\scanpage{55}
% archive foliation: 54
\[
F(C_\bullet^i)\unsure{-|\longrightarrow}\sigma(\rho(F))(C_\bullet^i)
=\operatorname{Hom}_k^\bullet(C_i^\bullet,\underbrace{\rho(F)}_{\Vert\;F(C_\bullet^\bullet)})
\simeq F(C_\bullet^i)
\]
\[
\begin{aligned}
F(\Psi_\bullet^i)&\longrightarrow\operatorname{Hom}_k^\bullet(\Psi_i^\bullet,F(C_\bullet^\bullet))
\simeq Z^i(F(C_\bullet^\bullet))\\
&\simeq\operatorname{Ker}(F(C_\bullet^i)\longrightarrow F(C_\bullet^{i+1}))\\
&\simeq F(\underbrace{\operatorname{Ker}(C_\bullet^i\unsure{-|\longrightarrow}C_\bullet^{i+1})}_{\Sigma\text{-Ker !}})
=F(\Psi_\bullet^i).
\end{aligned}
\]
Donc on voit que $u(L_\bullet):F(L_\bullet)\longrightarrow\sigma\rho(F)(L_\bullet)$
est un iso sur les $C_\bullet^i$ et les $\Psi_\bullet^i$ donc
sur \deleted{\unsure{toute}} la cat. additive $\unsure{\mathcal P|P}_0\supset K$ des sommes
de tels.) Pour avoir $\forall L_\bullet$, on \unsure{raisonne}
\inserted{\unsure{ess. par}}
par récurrence sur la longueur $n$ \deleted{\unsure{de $L_\bullet$}},
\inserted{de l'intervalle $[p,q]$ tel que $L_i=0$ si $i\notin[p,q]$.}
C'est vrai si $n=0$, ($L_\bullet$ somme $\simeq(\Psi_\bullet^i)^r$
\deleted{\unsure{$i>0$}}), on
suppose $n>0$ et suppose que c'est fait pour
$n-1$.
\diagram{\begin{tikzcd}[column sep=small,row sep=large]
&&\cdots\arrow[r]&0\arrow[r]&L_p\arrow[r,"\mathrm{id}"]&L_p\arrow[r]&0
&\left.\vphantom{\begin{matrix}L_p\\L_p\end{matrix}}\right\}
& (C_\bullet^{\unsure{p|p+1}})^r\;\deleted{\unsure{L_\bullet''}}\\
\cdots\arrow[r,dash]&0\arrow[r,dash]&L_{p+n}\arrow[r]&L_{p+n-1}\arrow[r,dashed]&L_{p+1}\arrow[u,"d_{p+1}"]\arrow[r]&L_p\arrow[u,"\mathrm{id}"]\arrow[r]&0\arrow[r]&0&L_\bullet\arrow[u]\\
&&&\underbrace{\hspace{5.8em}}_{\text{tronqué }L_\bullet'}&&&&&
\end{tikzcd}}
on trouve un \unsure{morph.} de suites exactes (où $L_p\simeq k^r$)
\diagram{\begin{tikzcd}[column sep=small,row sep=large]
0\arrow[r,dash]&(\Psi_\bullet^p)^r\arrow[r,dash]\arrow[d,equal]&L_\bullet\arrow[r]\arrow[d]&\deleted{\unsure{\Psi}}L_\bullet'\arrow[r]\arrow[d]&0\\
0\arrow[r]&(\Psi_\bullet^p)^r\arrow[r,dash]&(C_\bullet^p)^r\arrow[r,dash]&(\Psi_\bullet^{p+1})^r\arrow[r,dash]&0
\end{tikzcd}}
qui exprime $L_\bullet$ comme produit fibré
de $L_\bullet'$ et $(C_\bullet^p)^r$ sur $(\Psi_\bullet^{p+1})^r$ (i.e. comme
noyau d'un hom \inserted{$\Sigma$-\uline{épi}}
\[
(C_\bullet^p)^r\times L_\bullet^{\unsure{1|'}}\unsure{-|\longrightarrow}(\Psi_\bullet^{p+1})^r).
\]
Donc on conclut par

